\section{No-Other-Solutions Standard}
\label{sec:pr35-no-other-solutions-standard}

For each sector $s$, let \(\mathfrak A_s\) be the explicitly declared
admissible candidate set, let \(\sim_s\) identify gauge, basis, phase,
coordinate, algebraic, and unit conventions, and let \(F_s\) collect the
target-independent mathematical and closure constraints.  The meaningful
uniqueness claim requires
\begin{equation}
\boxed{
\{x\in\mathfrak A_s:F_s(x)=0\}/\!\sim_s
=
\{[x_{\rm PR}]\}.}
\label{eq:pr35-sector-uniqueness-standard}
\end{equation}
Verifying only \(F_s(x_{\rm PR})=0\) establishes existence, not uniqueness.
The proof must also show local isolation after unphysical null directions are
removed and exhaust every disconnected branch \emph{inside the same declared
class}.  A model in a broader theory belongs to a different theorem domain
unless it also satisfies all of \(F_s\).  Conversely, if two inequivalent
members satisfy every declared condition, the admissible quotient contains two
classes rather than the singleton required by
Eq.~\eqref{eq:pr35-sector-uniqueness-standard}.

\subsection{Equation Equivalence and Exhaustive Uniqueness}
\label{subsec:pr35-equation-equivalence}

Fix a sector $s$.  Let $\mathcal X_s$ be its configuration space after all
upstream fields, domains, normalizations, and identities have been fixed; let
$\mathcal Y_s$ be its residual space; and let $\mathcal G_s$ be the declared
group of gauge transformations, normalized basis changes, coordinate and unit
changes, and invertible field redefinitions that preserve the frozen ledger
and registered observable algebra.  A candidate equation is a covariant
residual map
\begin{equation}
 E:\mathcal X_s\longrightarrow\mathcal Y_s,
 \qquad
 Z(E):=\{x\in\mathcal X_s:E(x)=0\}.
 \label{eq:pr35-equation-zero-set}
\end{equation}

The structural candidate class $\mathfrak E_s^0$ is fixed before its equations
are solved and is saturated under $\sim_{\rm op}$, so an invertible rewriting
cannot be excluded merely by typography.  It is specified by field content, operator domain, locality or
finite-tier order, regularity, covariance, boundary grammar, representation
content, and inherited normalizations.  A target-independent closure ledger is
a map
\begin{equation}
 \Delta_s:\mathfrak E_s^0\longrightarrow[0,\infty)^{m_s}.
 \label{eq:pr35-closure-defect-map}
\end{equation}
An equation is admissible precisely when $\Delta_s(E)=0$.  Neither
$\mathfrak E_s^0$ nor $\Delta_s$ may contain a measured target, diagnostic
residual, fitted tolerance, or the desired answer inserted as a coefficient.
Unit normalization, zero incidence, completeness, nonredundancy, covariance,
and fixed-order minimality are admissible only when stated independently of the
value obtained after the equations are solved. Three relations must be distinguished:
\begin{enumerate}
 \item $E\sim_{\rm op}E'$, called \emph{operator equivalence}, when there are
 $g\in\mathcal G_s$ and $A_x\in\operatorname{GL}(\mathcal Y_s)$ at every $x$
 such
 that
 \begin{equation}
  E'(gx)=A_xE(x)
  \qquad\text{for every }x\in\mathcal X_s.
  \label{eq:pr35-operator-equivalence}
 \end{equation}
 This includes algebraic rearrangement, multiplication by a nowhere-zero
 factor, and allowed invertible changes of variables.
 \item $E\sim_{\rm sol}E'$, called \emph{solution equivalence}, when an allowed
 transformation maps the complete zero set of one equation bijectively onto
 that of the other, $gZ(E)=Z(E')$.
 \item $E\sim_{\rm phys}E'$, called \emph{physical equivalence}, when their
 quotient solution spaces are isomorphic and the isomorphism preserves every
 registered observable, multiplicity, boundary label, and physical spectral
 datum:
 \begin{equation}
  Z(E)/\mathcal G_s\cong Z(E')/\mathcal G_s.
  \label{eq:pr35-physical-equation-equivalence}
 \end{equation}
\end{enumerate}
Operator equivalence implies solution equivalence, and solution equivalence
under an allowed ledger-preserving transformation implies physical
equivalence.  The converses need not hold.  In particular, matching one
reported number is not equivalence.  Equations must agree on the complete
admissible domain and preserve the physical operator, not merely one
diagnostic evaluation.  For example, $u=0$ and $u^3=0$ have the same real zero
set, but their linearizations at the zero have different rank and they are not
operator-equivalent through an invertible residual multiplier.

\begin{theorem}[Normal-form exhaustion and alternative-equation trichotomy]
\label{thm:pr35-normal-form-exhaustion}
Suppose a sector has a parameter space $\Theta_s$ and a normal-form map
$\nu_s:\mathfrak E_s^0\to\Theta_s$ with the following independently proved
properties:
\begin{enumerate}
 \item $\nu_s$ is complete:
 $\nu_s(E)=\nu_s(E')$ if and only if $E\sim_{\rm op}E'$;
 \item admissibility is decided in normal form by a target-independent
 certificate $d_s:\Theta_s\to[0,\infty)^{m_s}$,
 \begin{equation}
  \Delta_s(E)=0
  \quad\Longleftrightarrow\quad
  d_s(\nu_s(E))=0;
 \end{equation}
 \item the certificate has exactly one zero,
 $d_s^{-1}(0)=\{\theta_s^\star\}$.
\end{enumerate}
Then
\begin{equation}
 \boxed{
 \{E\in\mathfrak E_s^0:\Delta_s(E)=0\}/\!\sim_{\rm op}
 =\{[E_s^\star]_{\rm op}\},}
 \label{eq:pr35-equation-class-singleton}
\end{equation}
and every admissible equation induces the same physical quotient solution
space as $E_s^\star$.  If, in addition, the canonical equation has an isolated
physical solution,
\begin{equation}
 Z(E_s^\star)/\mathcal G_s=\{[x_s^\star]\},
 \label{eq:pr35-canonical-solution-isolation}
\end{equation}
then the sector has one physical solution class as well.  Every proposed
equation $\widetilde E$ has exactly one of three physical dispositions:
it is physically equivalent to $E_s^\star$ on the complete declared domain
and is not a second physical equation; it is physically inequivalent but lies
in $\mathfrak E_s^0$, in which case it has a nonzero defect; or it is physically
inequivalent and lies outside $\mathfrak E_s^0$, in which case it belongs to an
enlarged theory outside the theorem.  An inequivalent in-class equation
satisfying every declared condition enlarges the admissible quotient and must
therefore be included in the theorem statement.
\end{theorem}

\begin{proof}
Let $E$ be admissible.  The certificate property gives
$d_s(\nu_s(E))=0$, so the unique-zero property gives
$\nu_s(E)=\theta_s^\star$.  If $E_s^\star$ is any representative with that
normal form, completeness gives $E\sim_{\rm op}E_s^\star$.  Hence all
admissible equations occupy one operator-equivalence class.
Equation~\eqref{eq:pr35-operator-equivalence} maps their zero sets
bijectively and uses only transformations in $\mathcal G_s$, so their physical
quotient solution spaces agree.  If
Eq.~\eqref{eq:pr35-canonical-solution-isolation} also holds, each of those
spaces is the same singleton.  The three dispositions are the same argument
and its contrapositive for in-class equations.  For an outside proposal, a
full-domain physical-equivalence test places it in the first disposition if it
is merely another representation; otherwise its inequivalent operator content
places it in the enlarged-theory disposition.
\end{proof}

The theorem is noncircular: the candidate class, quotient group, normal-form
classification, and defect ledger are fixed without using
$\theta_s^\star$.  That value first appears as the computed zero of the
certificate.  The substantive proof obligation is therefore to classify every
in-class equation into a faithful normal form and then exhaust the complete
zero set of its defect.

Finite discrete candidates are treated by exhaustive enumeration. Polynomial
and rational constraints are solved by exact elimination or complete analytic
parameterization. Continuous numerical roots require interval isolation and
exclusion of the complement. Operator claims require self-adjoint domains and
certified spectral enclosures. Radiative corrections require a complete
invariant-operator basis and an explicit power-counting rule.

Post-generation diagnostic agreement is not a generative admissibility
constraint. If
\(\mathcal E_s\) denotes external data and \(\mathcal C_s\) the construction
map, the firewall condition is
\begin{equation}
\mathcal E_s\cap\operatorname{Dom}(\mathcal C_s)=\varnothing .
\end{equation}
No residual or standard-deviation score may select a branch, coefficient,
truncation order, numerical resolution, or acceptance tolerance.

Thus ``no other solution'' never means that no differently written equation
can return the same number.  It means that no inequivalent operator or branch
in the declared PR closure class satisfies the complete structural ledger.
The global theorem may combine direct postulate theorems and proved canonical
closures, provided their status is explicit and every included sector contains
one zero-defect equivalence class.
